\section{Introduction}
\label{sec:introduction}

The canonical path for connecting a .NET web application to a PostgreSQL database has historically involved a chain of manual steps: installing the database engine, creating a database and user, writing a connection string into \texttt{appsettings.json}, and remembering to exclude credentials from source control. Each step is individually simple; in aggregate they represent a setup tax that compounds across team members, CI pipelines, and deployment environments. When the application is containerised the tax increases: the developer must now also author Docker Compose files, manage container lifetimes, and ensure the application does not start before the database is ready.

.NET Aspire~9 eliminates this tax at the local development layer. An AppHost project --- a thin \texttt{.exe} that hosts a \texttt{DistributedApplication} --- becomes the single entry point for the entire development stack. It declares containers as typed resources, wires connection strings into dependent projects at startup, enforces readiness ordering via \texttt{WaitFor}, and exposes a browser-based dashboard with aggregated logs, metrics, and traces. Postgres passwords are generated randomly on first run and stored in .NET User Secrets; they are never written to \texttt{appsettings.json} or committed to source control.

This article documents every concrete step required to add that orchestration layer to an existing .NET~9 solution. Section~\ref{sec:background} reviews the architectural primitives that Aspire introduces. Section~\ref{sec:apphost} covers the creation and structure of the AppHost project. Section~\ref{sec:wiring} details the canonical \texttt{Program.cs} resource wiring. Section~\ref{sec:servicedefaults} introduces the ServiceDefaults shared library and its observability configuration. Section~\ref{sec:web} describes the changes required in the orchestrated web project. Section~\ref{sec:secrets} explains the User Secrets flow and common credential pitfalls. Section~\ref{sec:discussion} addresses limitations and production generalisation. Section~\ref{sec:conclusion} concludes.

Throughout, the HobbitGame solution --- a Blazor Server application where an LLM-driven hobbit wanders a map of the Shire and persists state in PostgreSQL --- serves as the worked example, providing concrete project names, database resource names, and code excerpts that can be adapted to any other solution.
